% Options for packages loaded elsewhere
\PassOptionsToPackage{unicode}{hyperref}
\PassOptionsToPackage{hyphens}{url}
%
\documentclass[
  10pt,
]{article}
\usepackage{amsmath,amssymb}
\usepackage{iftex}
\ifPDFTeX
  \usepackage[OT1]{fontenc}
  \usepackage[utf8]{inputenc}
\else % if luatex or xetex
  \usepackage{unicode-math} % this also loads fontspec
  \defaultfontfeatures{Scale=MatchLowercase}
  \defaultfontfeatures[\rmfamily]{Ligatures=TeX,Scale=1}
\fi
\ifPDFTeX\else
  % xetex/luatex font selection
\fi
% Use upquote if available, for straight quotes in verbatim environments
\IfFileExists{upquote.sty}{\usepackage{upquote}}{}
\IfFileExists{microtype.sty}{% use microtype if available
  \usepackage[]{microtype}
  \UseMicrotypeSet[protrusion]{basicmath} % disable protrusion for tt fonts
}{}
\makeatletter
\@ifundefined{KOMAClassName}{% if non-KOMA class
  \IfFileExists{parskip.sty}{%
    \usepackage{parskip}
  }{% else
    \setlength{\parindent}{0pt}
    \setlength{\parskip}{6pt plus 2pt minus 1pt}}
}{% if KOMA class
  \KOMAoptions{parskip=half}}
\makeatother
\usepackage{xcolor}
\usepackage[margin=.82in]{geometry}
\setlength{\emergencystretch}{3em} % prevent overfull lines
\providecommand{\tightlist}{%
  \setlength{\itemsep}{0pt}\setlength{\parskip}{0pt}}
\setcounter{secnumdepth}{-\maxdimen} % remove section numbering
\ifLuaTeX
  \usepackage{selnolig}  % disable illegal ligatures
\fi
\usepackage{bookmark}
\IfFileExists{xurl.sty}{\usepackage{xurl}}{} % add URL line breaks if available
\urlstyle{same}
\hypersetup{
  pdftitle={Path variation and fractional planks},
  pdfauthor={Leo Gao, AI-assisted research},
  hidelinks,
  pdfcreator={LaTeX via pandoc}}

\title{Path variation and fractional planks}
\author{Leo Gao, AI-assisted research}
\date{}

\renewcommand{\labelitemi}{$\bullet$}
\begin{document}
\maketitle

\section{A path-variation bound for fractional planks with two equal
normal
coefficients}\label{a-path-variation-bound-for-fractional-planks-with-two-equal-normal-coefficients}



\textbf{Theorem.} Let \(K\) be a convex subset of the unit three-cube
meeting all six facets, and suppose \(K\) lies in \(x+y+az=t\) with
\(a\ge0\). Then its fractional plank covering number satisfies \[
\tau^*(K)\ge \max\{1-a,0\}.
\] In particular the conjectured threshold \(2/3\) holds for \textbf{all
offsets when \(0\le a\le1/3\)}. The previous convolution bound proved it
only for \(a\le5-2\sqrt6\approx0.10102\). The construction below uses no
convolution and does not require a single inscribed triangle. This does
not resolve the whole two-equal-normal problem.

A fractional cover consists of nonnegative measurable profiles \(f_i\)
of the three coordinates with \(\sum_i f_i(x_i)\ge1\) on \(K\); the cost
is \(\sum_i\int_0^1 f_i\). A finite nonnegative measure on \(K\) with
each marginal dominated by Lebesgue measure certifies cost at least its
mass. This is the dual viewpoint of
\href{https://doi.org/10.1007/s00454-001-0088-x}{Aharoni--Holzman--Krivelevich--Meshulam,
Fractional Planks (2002)}.

\subsection{1. Trimming a shear by downward path
variation}\label{trimming-a-shear-by-downward-path-variation}

For a continuous piecewise affine function \(g:[0,1]\to[0,1]\), let
\(V_-(g)=\int_0^1(-g'(s))_+\,ds\), the sum of the decreases along its
affine pieces. Let \(\nu\) denote the distribution of \[
F(s)=s+a g(s)+b
\] for uniform \(s\in[0,1]\), with \(a\ge0\) and any real \(b\). For
every Borel set \(E\subseteq\mathbb R\) of finite length, \[
\nu(E)-|E|\le aV_-(g). \tag{1}
\] To see this, take \(H(u)=\int_0^u\mathbf 1_E(v)\,dv\), interpreted as
a signed integral for \(u<0\). Direct substitution on each affine piece
of \(F\) gives \[
\int_0^1 F'(s)\mathbf 1_E(F(s))\,ds
 =H(F(1))-H(F(0))\le |E|.
\] The identity also holds on a constant piece, where both sides of that
piece's identity are zero. Since \(1=F'(s)-ag'(s)\) almost everywhere,
(1) follows. This argument handles folds and atoms without treating them
separately.

The same bound holds for a mixture of such paths, with its right-hand
side averaged over the path weights. In particular, if each path has
\(V_-(g)\le1\), the projected shear law \(\nu\) has nonnegative excess
over Lebesgue measure of total mass at most \(a\). Explicitly, its
decomposition \(\nu=q\,dv+\nu_s\) gives \[
\sup_E(\nu(E)-|E|)=\int(q-1)_+\,dv+\nu_s(\mathbb R).
\] Here the laws in use have compact support, so the supremum can be
taken over finite-length sets. To trim a measure with this shear
marginal, discard its singular shear part and on the absolutely
continuous part keep the fraction \(\min(1,1/q(v))\) at shear value
\(v\) (take the fraction one when \(q(v)=0\)). The shear marginal
becomes \(\min(q,1)\,dv\). The retained mass is at least \(1-a\) for an
original probability, and no other marginal has increased.

\subsection{2. An ordered square contact hull supports bounded-variation
paths}\label{an-ordered-square-contact-hull-supports-bounded-variation-paths}

Let a convex set in the \((s,z)\) unit square contain four contacts \[
L=(0,C),\quad R=(1,D),\quad P=(A,0),\quad Q=(B,1),
\qquad A\ge B,\quad C\ge D.
\] There is a probability in their hull with both \(s,z\) uniform, which
can be built from paths of downward \(z\)-variation at most one. First
assume \(0<B<A<1\) and \(0<D<C<1\). Consider the following four
piecewise affine paths, each parameterized by its first coordinate from
zero to one: \[
\begin{array}{c|c|c}
\text{path}&\text{vertices in order}& V_-\text{ in }z\\ \hline
\mathrm{lo}& L,P,R & C\\
\mathrm{up}& L,Q,R & 1-D\\
\mathrm{vert}& L,Q,P,R & 1\\
\mathrm{diag}& L,R & C-D
\end{array}
\] Every edge stays in the hull by convexity. Give each path the law of
uniform \(s\); all four first marginals are automatically uniform. Their
\(z\) densities are constant on the three open intervals
\((0,D),(D,C),(C,1)\). To find a mixture with uniform \(z\), associate
to a path the vector \[
(\text{density on }(0,D)-\text{density on }(D,C),\
 \text{density on }(C,1)-\text{density on }(D,C)).
\] From the elementary density formula for a segment, the respective
four vectors are \[
\begin{array}{c|cc}
\mathrm{lo} & (1-A)/D &-A/C\\
\mathrm{up} & -(1-B)/(1-D)& B/(1-C)\\
\mathrm{vert}&(1-A)/D & B/(1-C)\\
\mathrm{diag}&-1/(C-D)&-1/(C-D).
\end{array}
\] For example, the middle edge \(QP\) contributes the same density
\(A-B\) on the three intervals of the vertical path, and cancels from
its two differences.

There is one vector in each open quadrant, so the origin belongs to
their convex hull. One direct check is that for every nonzero linear
functional, choose the displayed vector whose coordinate signs oppose
its nonzero coefficients: its value is negative, contradicting any
separating functional. A mixture giving the zero vector has equal
density on the three intervals, hence density one since it has total
mass one. Each of its paths has downward variation at most one, as
required. This uses the same kinds of boundary segments and diagonals as
\href{https://msp.org/pjm/1988/135-2/pjm-v135-n2-p06-s.pdf}{Gardner's
Theorem 1, case (ii)}; the routing into full-width paths and the shear
estimate are the additional steps.

For degenerate ordered contacts, replace each of \(A,B,C,D\) by its
convex interpolation with fixed strictly ordered interior values. The
convex hulls of the adjusted contacts converge to the original hull.
Apply the preceding construction, then trim as in Section 1 for any
fixed \(a,b\). The measures are supported in a fixed compact set in the
\((s,z)\) plane, so a subsequence converges weakly. Its limit is
supported on the actual contact hull. All three marginal dominations,
including the shear marginal, survive the limit by testing with
nonnegative continuous functions; their masses converge and each is at
least \(1-a\). This form of the conclusion is all that is needed; no
claim about the piecewise-path form of the limiting untrimmed measure is
necessary.

\subsection{3. Apply the ordering to the two equal
coefficients}\label{apply-the-ordering-to-the-two-equal-coefficients}

We may restrict to the compact convex hull of finitely many facet
contacts. When \(a=0\), the plane is \(x+y=1\), and Gardner's
two-coordinate theorem applied to \((x,z)\) gives full uniform marginals
in all three coordinates after lifting. Now let \(a>0\). Packedness
implies \(1\le t\le1+a\) (the stronger feasible inequalities are
unnecessary here). Choose one point of \(K\) on each of \(z=0,1\). If
their respective \((x,y)\) coordinates are \((A_x,A_y)\) and
\((B_x,B_y)\), then \[
 (A_x-B_x)+(A_y-B_y)=a>0.
\] At least one of the two horizontal coordinates, call it \(s\),
satisfies \(A_s\ge B_s\). In the \((s,z)\) projection choose contacts
\((0,C),(1,D)\). The unselected horizontal coordinate lies in \([0,1]\),
so \[
 C\ge (t-1)/a\ge D.
\] Indeed at \(s=0\) the other value is \(t-aC\le1\), and at \(s=1\) the
other value is \(t-1-aD\ge0\). Thus the four chosen contacts satisfy
Section 2.

Lift the paths back to the plane. The remaining horizontal coordinate is
\(t-s-az\), whose reflection is \[
1-t+s+az .
\] This is exactly the shear of Section 1 with \(b=1-t\). After
trimming, the projected \(s,z\) marginals and the reflected remaining
marginal are all dominated by Lebesgue measure. Reflection preserves
this domination for the original coordinate, so the measure is a
matching in \(K\) of mass at least \(1-a\). The degenerate-case limiting
argument applies even if the auxiliary strict-order hulls are not
contained in the original plane section of the cube: they are used only
in the projection before taking the limit back to the exact contacts.

The resulting bound proves the new all-offset range. The argument does
not by itself establish the sharp threshold for \(a>1/3\); its loss
bound is an upper bound on the mass discarded, not a counterexample in
the remaining parameter range.

\section{Strictness of the path-variation
bound}\label{strictness-of-the-path-variation-bound}



\textbf{Theorem.} For every \(0<a<1\), every packed convex \(K\) in
\(x+y+az=t\) has \(\tau^*(K)>1-a\). On each fixed compact interval of
coefficients \(J\subset(0,1)\) this improvement over \(1-a\) is uniform
over all feasible offsets and all packed \(K\). In particular there is
some nonexplicit \(\epsilon>0\) for which the sharp conjectured bound
\(2/3\) holds at \textbf{every} feasible offset for
\(0\le a\le 1/3+\epsilon\).

Use the previous note's notation for a chosen horizontal coordinate
\(s\), the other reflected coordinate \(W=1-t+s+az\), and the contact
hull with vertices \(L=(0,C), R=(1,D), P=(A,0), Q=(B,1)\) in \((s,z)\).
Choose \(s\) so that \(A>B\), which is possible because the sum of the
two candidate differences is \(a>0\). Put \(c=(t-1)/a\). For interior
offsets \(1<t<1+a\), we have \(0<c<1\) and \(C\ge c\ge D\).

\subsection{1. Interior offsets: a strictly better path mixture, or an
augmentable
antidiagonal}\label{interior-offsets-a-strictly-better-path-mixture-or-an-augmentable-antidiagonal}

If \(0<D<C<1\) and \(0<B<A<1\), the four density-difference vectors in
the main proof have the origin in the interior of their hull (every
nonzero functional takes both signs). They therefore admit a solution
assigning positive weights to all four paths: first allocate a
sufficiently small common positive weight to each, then represent the
small residual displacement using their convex hull. Paths other than
the vertical route have \(V_-<1\). Hence the average variation is less
than one and clipping the shear marginal loses strictly less than \(a\).

If still \(0<D<C<1\) but \(A=1,B>0\), the lower and the vertical path
vectors are respectively \((0,-1/C)\) and \((0,B/(1-C))\); their
positive mixture gives the zero vector, and the lower path has
\(V_-=C<1\). If \(B=0,A<1\), instead use the upper and the vertical
paths, whose first components are opposite and whose second components
vanish; \(V_-\) for the upper path is \(1-D<1\). In such cases a
vertical segment with zero \(s\) width at a path endpoint is simply
omitted when parameterizing by \(s\). If \(A=1\) and \(B=0\), the
segment from \(Q\) to \(P\) is a full antidiagonal.

If \(C=D=c\in(0,1)\), the lower path carries a \(z\) marginal uniform on
\([0,c]\), the upper path one on \([c,1]\). Mix with weights \(c,1-c\)
to get a uniform \(z\) marginal and average downward variation
\(c^2+(1-c)^2<1\), including when one of these paths skips a zero-width
endpoint segment.

If \(C=1\) and \(0<D<1\), the lower path and the straight diagonal
\(LR\) suffice when \(A<1\). The lower path density on \((0,D)\) is
\(A+(1-A)/D>1\), and both paths have constant density on each of
\((0,D),(D,1)\). Assign the lower path the reciprocal of that density
and the straight path the remaining positive mass. The density is then
one on both intervals by normalization, and the straight path has
\(V_-=1-D<1\). If \(A=1\) the segment \(LP\) is instead a full
antidiagonal. In the case \(0<C<1,D=0\), use the upper path and \(LR\)
when \(B>0\): the upper path has density \((1-B)+B/(1-C)>1\) on
\((C,1)\), so it receives the reciprocal of that density, and the rest
of the mass uses \(LR\) with \(V_-=C<1\). For \(B=0\) the segment \(QR\)
is a full antidiagonal. The last possibility \(C=1,D=0\) gives a full
antidiagonal directly via \(LR\).

Whenever there is a full antidiagonal in \((s,z)\), its image under
\(W\) is an interval of length \(1-a\). Put mass \(1-a\) uniformly on
the antidiagonal. Its \(s,z\) marginals have density \(1-a<1\) and its
\(W\) marginal has density one on that interval only. Since \(K\) is
packed in the remaining coordinate, it contains a point whose \(W\) lies
strictly outside the interval (at least one of 0 and 1 does). Connect it
to an antidiagonal point that differs from it in all three projected
coordinates. This is possible: the antidiagonal is injective in \(s,z\),
so at most two of its points are excluded, and every point on it differs
in \(W\). A short initial subsegment lies entirely outside the saturated
\(W\) interval and has nonzero length in all projections. Give it a
sufficiently small positive uniform mass. It fits under the \(s,z\)
density slack \(a\) and the \(W\) density cap one, proving the strict
improvement.

\subsection{2. Endpoint offsets}\label{endpoint-offsets}

At \(t=1\), the two cube vertices \((1,0,0)\) and \((0,1,0)\) are forced
into \(K\) by the corresponding facet contacts. Choose any \(z=1\)
contact. Its \(s\) coordinate \(B\) is in \([0,1-a]\), and the
projection contains the triangle \((0,0),(1,0),(B,1)\) in \((s,z)\).
Here \(W=s+az\). Consider the route from \((0,0)\) to \((B,1)\) to
\((1,0)\) with uniform \(s\) parameter. If \(B=0\) omit the first edge
of zero \(s\) width. The \(z\) marginal is uniform, since the two edges
contribute respective constant densities \(B\) and \(1-B\) on \((0,1)\).

The first edge projects to \(W\in[0,B+a]\) with density \(B/(B+a)\). The
second projects to \([B+a,1]\); retain the fraction \((1-B-a)/(1-B)\) of
its mass, making its \(W\) density one. If \(B+a=1\) its \(W\)
projection is constant and the retained fraction is zero. The resulting
measure has mass \(1-a\). Its \(z\) marginal has density \(1-a\)
throughout \((0,1)\); its \(s\) marginal has positive slack on
\((B,1)\); and its \(W\) marginal has positive slack on \((0,B+a)\).

Inside this triangle there is an open region with \(s>B\) and \(W<B+a\).
For example, start at any baseline \(s\in(B,B+a)\) and raise \(z\) a
sufficiently small positive amount. Choose a short segment in this
region varying all three coordinates, and give it a sufficiently small
positive uniform mass, respecting the three slacks. This yields a strict
improvement at \(t=1\). Reflect all coordinates for \(t=1+a\).

\subsection{3. Uniformity on a compact coefficient
range}\label{uniformity-on-a-compact-coefficient-range}

Suppose there were no uniform improvement on a compact
\(J\subset(0,1)\). For a sequence violating such an improvement, choose
six contacts and pass to a subsequence on which their coefficients,
offsets, and all contact points converge. The equal-normal pair of first
two coefficients is retained; call the limiting third coefficient
\(a_*>0\). The limiting contact hull has interior in its plane. Indeed,
a packed segment would have absolute coordinate increments all equal to
one, which cannot solve \(\pm1\pm1\pm a_*=0\) for \(a_*\in(0,1)\). The
strict result gives a dominated measure of mass \(m>1-a_*\) in the
limiting hull.

Contract slightly toward an interior point and scale its mass by the
contraction factor, preserving all marginal caps while keeping mass
above \(1-a_*\). Convolve in the plane with a smooth compactly supported
nonnegative probability kernel of sufficiently small support. This
preserves every cap by convolution and puts a smooth density compactly
inside the limiting hull. Use \((s,z)\) coordinates, with the same
choice of \(s\) for all hulls in this argument. The nearby contact hull
projections converge in Hausdorff distance, so they eventually contain
this compact support (use the positive support-function margin). Its
\(s,z\) densities remain unchanged; the third density for the linear
form \(t-s-az\) varies uniformly with \(a,t\) by the integral formula
for projecting a smooth compactly supported density. Divide by
\(1+o(1)\) to restore all three caps, obtaining a uniform mass excess
above \(1-a_n\) for large \(n\), a contradiction. Apply this around
\(a_*=1/3\) for the final consequence.

\end{document}
